\begin{ExerciseBlock}
\ExerciseHeading{01}{Drehimpulsvektor}{angular_momentum_angular_momentum_vector}
Ein Körper der Masse $m = 10\,$kg bewegt sich um den Koordinatenursprung $(0,0,0)$. Momentan hat der Körper die Koordinaten $\vec{r}$ und die Geschwindigkeit $\vec{v}$.

\[
  \vec r=\begin{pmatrix}2\\3\\-1\end{pmatrix}\,\unit{\metre}
  \qquad\text{und}\qquad
  \vec v=\begin{pmatrix}1\\0\\3\end{pmatrix}\,\unit{\metre\per\second}.
\]

\begin{enumerate}
	\item Wie kann man anhand der Angaben erkennen, dass der Körper sich nicht auf einer Kreisbahn um den Koordinatenursprung bewegt?

	\item Berechnen Sie den Drehimpulsvektor $\vec{L}$ des Körpers.

	\item Zeigen Sie, dass $\vec{L} \perp \vec{r}$ und $\vec{L} \perp \vec{v}$.

\end{enumerate}
\end{ExerciseBlock}
\UseExerciseSolution{angular_momentum_angular_momentum_vector}
